\section{Pruebas funcionales y resultados}\label{sec:pruebas-funcionales}

El 16 de octubre de 2026 se auditó el conjunto de pruebas de MiSendero y se amplió su alcance para obtener evidencia instrumentada de cobertura. Se distinguieron dos indicadores: la proporción de casos aprobados y la proporción de código ejecutado por esos casos. La aprobación de 93 pruebas, declarada en una versión anterior, no permite inferir una cobertura del 85\,\%. Tampoco corresponde trasladar ese recuento a otra revisión del software: al comenzar esta auditoría, la revisión \texttt{c33b658b2556} contenía 185 casos del cliente y 93 del backend, todos aprobados antes de incorporar las nuevas pruebas.

\subsection{Método, indicador y alcance}

Se adoptó la cobertura de líneas ejecutables indicada en RNF-11 y en la matriz interna de medición. No se dispuso de una devolución original de la cátedra que precisara el indicador; por ello, se recomienda confirmar formalmente si el umbral corresponde a líneas, sentencias, ramas u otra medida. La evidencia presentada acredita el criterio de líneas del alcance declarado y no permite sustituirlo por otro indicador.

La instrumentación Babel/Istanbul contabilizó el código propio JavaScript y TypeScript del cliente y del backend, incluidos los archivos no importados por las pruebas. El inventario contiene 141 archivos del cliente y 31 del backend. Comprende rutas de entrada, pantallas, componentes, hooks, servicios, persistencia, lógica de inferencia, configuración de distribución, scripts auxiliares, API y código de la página web de gestión de cuentas. Las declaraciones sin instrucciones ejecutables se conservan en el inventario y aportan cero líneas al denominador.

Se separaron el código de pruebas y sus mocks, las dependencias de terceros, los archivos generados y temporales, los documentos, los recursos gráficos, las migraciones SQL y los datos JSON. Estos últimos pueden intervenir en las pruebas, pero no constituyen instrucciones JavaScript o TypeScript instrumentables. El código Python de investigación se midió íntegramente con otra herramienta y se informa de manera independiente. No se retiraron pantallas o servicios por presentar baja cobertura ni se incorporaron directivas para ignorar código relevante.

La cobertura conjunta se calculó mediante la suma de contadores de líneas de dos inventarios disjuntos:
\[
C_{\mathrm{líneas}}=100\times
\frac{\sum_i L_{i,\mathrm{cubiertas}}}{\sum_i L_{i,\mathrm{ejecutables}}}.
\]
No se promediaron porcentajes de archivos o componentes. Ambos conjuntos utilizan el mismo proveedor de instrumentación y expresan contadores de código fuente. Las ramas, las sentencias y las funciones se conservaron como indicadores adicionales, con denominadores propios.

\subsection{Herramientas y procedimiento reproducible}

El cliente utiliza Jest 29.7.0 con jest-expo 57.0.5, React Test Renderer y mocks de los límites nativos. Se ejecutaron pantallas y componentes reales, geometrías de Three.js y operaciones SQL mediante el adaptador de pruebas de SQLite, basado en \texttt{node:sqlite}. El backend utiliza Jest 30.4.2 y Supertest sobre Express, con PostgreSQL emulado mediante pg-mem 3.0.14. Estos entornos verifican contratos y comportamiento de la lógica implementada; no equivalen a ejecutar los binarios móviles ni un servidor PostgreSQL de producción.

La ejecución final se realizó en Windows, con Node.js 24.19.0 y npm 11.17.0. Se detectó que las dependencias inicialmente instaladas correspondían a versiones anteriores de Expo y React. Se comprobó una instalación limpia del lockfile, sin cambiar sus versiones, mediante el siguiente procedimiento desde la raíz del repositorio:

\begin{verbatim}
npm ci --legacy-peer-deps --ignore-scripts
npm ci --prefix backend --ignore-scripts
npm run typecheck
npm run test:coverage -- --output=coverage/final 
    --randomize --seed=20261016 --showSeed
\end{verbatim}

La última orden debe escribirse en una sola línea. La opción \texttt{--legacy-peer-deps} reproduce el grafo bloqueado sin resolver nuevamente los pares nativos opcionales de Expo. La instalación ordinaria sin esa opción falló; esta condición queda documentada y no se interpreta como validación de compatibilidad de una compilación nativa.

El comando de cobertura fija la zona horaria UTC, ejecuta ambos conjuntos y verifica el inventario. Devuelve un estado de error ante pruebas fallidas u omitidas, archivos ausentes, discrepancias entre dependencias instaladas y bloqueadas o una cobertura conjunta de líneas inferior al 85\,\%. Dos pruebas adicionales del agregador verifican la suma ponderada, la detección de archivos no importados y el umbral sin redondear. La integración continua utiliza el mismo procedimiento y conserva los reportes incluso ante fallos.

\subsection{Comportamientos evaluados}

La ampliación incorporó 144 casos sobre comportamientos observables, sin debilitar las expectativas existentes. La Tabla~\ref{tab:technical-validation} sintetiza los principales escenarios.

\begingroup
\small
\setlength{\tabcolsep}{4pt}
\renewcommand{\arraystretch}{1.12}
\begin{longtable}{|p{3.1cm}|p{11.1cm}|}
\caption{Escenarios críticos incluidos en la auditoría automatizada}\label{tab:technical-validation} \\
\hline
\textbf{Componente} & \textbf{Comportamientos comprobados} \\
\hline
\endfirsthead
\multicolumn{2}{r}{\footnotesize Continuación de la tabla} \\
\hline
\textbf{Componente} & \textbf{Comportamientos comprobados} \\
\hline
\endhead
\hline
\endfoot
\hline
\endlastfoot
Operación offline & Consulta del contexto local; progreso y cancelación de descarga; rechazo de paquetes incompletos; verificación de archivos antes de abrir; persistencia de reportes cuando falla la API. \\
\hline
Sincronización & Aislamiento por cuenta, reintentos, idempotencia, reanudación por conectividad o primer plano y distinción entre operación pendiente y confirmada. \\
\hline
Ubicación y recorrido & Denegación de permisos; rechazo de puntos antiguos, imprecisos, fuera de orden o con velocidad inverosímil; distancia persistida; pausa y reanudación; aviso de desvío y selección de coordenadas próximas al trazado. \\
\hline
Persistencia y acceso & Sesiones y reseñas en SQLite; caché de planes por usuario; registro y autenticación; aceptación de términos independiente de otras decisiones; revocación, baja de cuenta y prevención de reenvíos pendientes. \\
\hline
Asistencia y emergencias & Cola de SOS, estados de planes, confirmación de regreso y reintentos del adaptador de notificaciones con la misma clave de idempotencia. La aceptación HTTP no se representa como entrega al contacto. Los datos simulados mantienen el SOS en modo de prueba. \\
\hline
Mapas, relieve y métricas & Capas locales y gestos; geometría y liberación de recursos de Three.js; ausencia explícita de relieve sin DEM; perfil de elevación accesible; contratos de HealthKit y Health Connect; consentimiento previo a la lectura y separación de muestras simuladas. \\
\hline
Gestión web y datos externos & Exportación y corrección de datos, confirmación de eliminación, moderación y representación de contenido como texto; rechazo de fotografías sin metadatos de licencia; errores de altimetría y pronóstico. \\
\hline
\end{longtable}
\endgroup

La ejecución con orden aleatorio detectó siete fallos en casos existentes: uno del cliente y seis del backend. Sus causas fueron precondiciones compartidas o dependencia de registros creados por otro caso. Se aislaron los datos de usuarios, actividades, senderos y cola; cuando un caso necesitaba un antecedente, pasó a crearlo explícitamente. Se conservaron las expectativas originales. El límite de intentos de autenticación se amplió únicamente en la configuración de pruebas para que la creación de fixtures no agotara un presupuesto compartido.

\subsection{Resultados instrumentados}

Con el inventario completo y los casos iniciales, la cobertura conjunta fue 2744/6123 líneas, equivalente al 44{,}81\,\%. La medición anterior limitada a archivos importados ofrecía porcentajes superiores, pero no acreditaba el universo completo. Tras ampliar las pruebas, el denominador permaneció en 6123 líneas y se cubrieron 5357: 2613 líneas adicionales.

\begin{table}[H]
\centering
\small
\setlength{\tabcolsep}{4pt}
\renewcommand{\arraystretch}{1.15}
\begin{tabular}{|p{4.6cm}|p{2cm}|p{2cm}|p{4cm}|}
\hline
\textbf{Conjunto} & \textbf{Suites} & \textbf{Casos} & \textbf{Resultado} \\
\hline
Cliente y servicios móviles & 44 & 308 & Aprobados; sin omisiones. \\
\hline
Backend, API y script web & 16 & 114 & Aprobados; sin omisiones. \\
\hline
\textbf{Total del producto JS/TS} & \textbf{60} & \textbf{422} & \textbf{Sin fallos.} \\
\hline
\end{tabular}
\caption{Pruebas del producto ejecutadas el 16 de octubre de 2026.}
\label{tab:resultados-pruebas}
\end{table}

\begin{table}[H]
\centering
\small
\setlength{\tabcolsep}{4pt}
\renewcommand{\arraystretch}{1.15}
\begin{tabular}{|p{2.7cm}|p{3.1cm}|p{2cm}|p{2cm}|p{2cm}|}
\hline
\textbf{Componente} & \textbf{Líneas cubiertas / total} & \textbf{Sentencias} & \textbf{Ramas} & \textbf{Funciones} \\
\hline
Cliente & 4539/5190 (87{,}45\,\%) & 83{,}67\,\% & 73{,}11\,\% & 77{,}01\,\% \\
\hline
Backend & 818/933 (87{,}67\,\%) & 85{,}50\,\% & 76{,}28\,\% & 81{,}92\,\% \\
\hline
\textbf{Conjunto} & \textbf{5357/6123 (87{,}49\,\%)} & 83{,}96\,\% & 73{,}42\,\% & 77{,}56\,\% \\
\hline
\end{tabular}
\caption{Cobertura de código fuente con inventario completo y agregación ponderada.}
\label{tab:cobertura-instrumentada}
\end{table}

La medición supera el mínimo del 85\,\% de líneas tanto por componente como en conjunto. No alcanza ese umbral conjunto en sentencias, ramas o funciones. La aprobación del 100\,\% de los casos de la Tabla~\ref{tab:resultados-pruebas} se interpreta de manera independiente de estos porcentajes.

Los reportes verificables se conservan en \path{docs/coverage-audit-2026-10-16/}, junto con el manifiesto SHA-256 de código, pruebas, configuración y evidencias. Se entregan JSON de contadores y resultados, CSV por archivo, LCOV, HTML con detalle de líneas y un archivo comprimido con las ejecuciones. Los informes se regeneran en \path{coverage/final/}. El registro de entorno identifica las versiones realmente instaladas y sus lockfiles.

Una segunda ejecución utilizó la semilla 73 y dos procesos, frente a la semilla 20261016 y ejecución secuencial de la medición principal. Ambos órdenes aprobaron los mismos casos y reprodujeron los cuatro contadores agregados. La instalación limpia también reprodujo la cobertura. Estas comprobaciones aportan evidencia de independencia y repetibilidad de los escenarios ejecutados; no constituyen una demostración de ausencia de todo comportamiento no determinista posible.

\subsection{Brechas y límites de la evidencia}

Permanecen 766 líneas del producto sin cubrir. Entre las brechas se encuentran el arranque y el enrutamiento inicial, la pantalla específica de recuperación ante desvíos, ramas de actualización y eliminación de copias offline, fallos adicionales de almacenamiento y algunos scripts de mantenimiento. Estos archivos continúan contabilizados. Se recomienda priorizar las ramas de error y los límites de permisos y conectividad por encima de casos que únicamente importen módulos o comprueben su existencia.

Las 14 pruebas Python de investigación se ejecutaron con unittest y coverage.py 7.6.1, sobre todos los módulos de \path{training/python/}, incluidos los no importados y excluyendo solamente sus pruebas. Alcanzaron 485/696 líneas (69{,}68\,\%) y 196/301 ramas (65{,}12\,\%); el indicador combinado de coverage.py fue 68{,}30\,\%. Este conjunto no alcanza el 85\,\% y se mantiene separado por su finalidad de investigación y por las diferencias de instrumentación. Sus resultados tampoco acreditan precisión predictiva o habilitación de modelos en campo.

La cobertura de código no sustituye las pruebas de extremo a extremo, las pruebas de dispositivos ni la evaluación de usabilidad. Los límites nativos, la red y los proveedores se controlaron mediante dobles de prueba. Continúan requiriendo evidencia independiente los permisos y la revocación efectivos en iOS y Android, el GPS y la batería en recorridos reales, la continuidad en segundo plano, la representación GPU, HealthKit y Health Connect en hardware físico, las migraciones sobre PostgreSQL real, la entrega de mensajes a contactos y la respuesta de servicios externos. Tampoco se efectuó en esta auditoría una nueva evaluación de usabilidad.

